\item \points{4b}

For Q-learning in continuous states, we need to use function approximation. The first step of function approximation is extracting 
features from the given state. Feature extractors of different complexities work well with different problems: linear and polynomial feature 
extractors that work well with simpler problems may not be suitable for other problems. For the mountain car task, we are going to use 
a Fourier feature extractor. 
    
For the Fourier feature extractor, the goal is to approximate Q-value using a sum of sinusoidal components: for state $s = [s_1, \ldots, s_k]$ and maximum coefficient $m$, the feature extractor $\phi$ is:
$$\phi(s, m) = [\cos(0), \cos(s_1\pi), \ldots, \cos(s_k\pi), \cos(2s_1\pi), \cos((s_1+s_2)\pi), \ldots, \cos((ms_1 + ms_2 + \ldots + ms_k)\pi)]$$

Note that $\phi(s, m) \in \mathbb{R}^{(m+1)^k}$. For those interested, look up Fourier approximation and how effective Fourier features are in different settings.

In the code you implement, there is another variable \texttt{scale} which is applied to different dimensions of the input state. We give you an example of how to compute this Fourier feature with $s = (s_1, s_2)$, \texttt{scale}$=[d_1, d_2]$ and $m=2$. We find all the entries by considering all the combinations of scaled entries of the state as below.
\begin{table}[h!]
    \centering
    \lstDeleteShortInline|                                                
    \begin{tabular}{|c|c|c|c|}
        \hline
        & 0 & $d_1 s_1$ & $2 d_1 s_1$ \\
        \hline
        0 & 0 & $d_1 s_1$ & $2 d_1 s_1$ \\
        \hline
        $d_2 s_2$ & $d_2 s_2$ & $d_1 s_1 + d_2 s_2$ & $2 d_1 s_1 + d_2 s_2$ \\
        \hline
        $2 d_2 s_2$ & $2 d_2 s_2$ & $d_1 s_1 + 2 d_2 s_2$ & $2 d_1 s_1 + 2 d_2 s_2$ \\
        \hline
    \end{tabular}
    \lstMakeShortInline|
\end{table}

From this, our feature extractor will return a numpy array with 9 elements - \\
$[\cos 0, \cos d_1s_1\pi, \cos 2d_1s_1\pi, \cos d_2s_2\pi, \cos (d_1s_1 + d_2s_2)\pi, \dots, \cos (2d_1s_1 + 2d_2s_2)\pi]$. If there is a third dimension $s_3$, we will repeat the "outer-sum" with the nine entries in the table and $(0, d_3s_3, 2d_3s_3)$, before multiplying by $\pi$ and applying cosine to the resulting 27 elements.

Implement \texttt{fourier\_feature\_extractor} in \texttt{submission.py}. 

Looking at \texttt{util.polynomialFeatureExtractor} may be useful
for familiarizing oneself with numpy.

\expect{An implementation of \texttt{fourier\_feature\_extractor} in \texttt{submission.py}}


